\section{Completing the hardness proof}\label{sec:completion}

The parameter order is
\[
 (t,\eps,\mu)\longrightarrow(s,m,p_*,\theta)
 \longrightarrow b\longrightarrow n\longrightarrow T\longrightarrow a.
\]
The extraction constants fix the required game soundness $b$ before the
alphabet dimension $n$ is known. Repetition is then chosen to hide the
planted coordinate, and the game's completeness error $a$ is chosen last.

\begin{proof}[Proof of \Cref{thm:gap}]
Fix $t$ and $\eps$ as in the theorem. Take $\delta=\eps$. This satisfies $0<\delta<1-B(t)$ because
$\eps<(t-B(t))/4<(1-B(t))/4$. Choose a
rational $0<\mu<1/2$ with $t\mu\le\eps/4$.
Apply \Cref{prop:extraction} with $t,\delta,\mu$, fixing
$s,m,p_*,\theta$, and hence $k,d,r$. Choose a rational
\[
 0<b<2^{-2s}\frac{p_*\theta}{2}.
\]
By \Cref{prop:affine-games}, this fixes an alphabet dimension $n$ for
which the game has arbitrarily small completeness error. Choose an
integer $T$ satisfying
\[
 T\ge\frac{2^{n(2k-2s)}}{(p_*\theta)^2}.
\]
Finally choose a positive rational $a<1-b$ small enough that
$2Ta\le\eps/4$, and start with the NP-hard game gap
$\Val(\mathcal G)\ge1-a$ versus $\Val(\mathcal G)\le b$.

In the YES case, \Cref{lem:completeness} gives cut value at least
$(1+t)/2-\eps/2$, which is stronger than required.
In the NO case, suppose a cut had value at least
$(1+B(t))/2+\eps/2$. By \Cref{lem:cut-stability}, its odd average has
expected slice stability at least $B(t)+\eps$.
\Cref{prop:extraction,prop:decoding} would then imply
$\Val(\mathcal G)\ge2^{-2s}p_*\theta/2>b$, a contradiction.
Thus every cut has value less than $(1+B(t))/2+\eps/2$.

For completeness, the test is a deterministic polynomial-size graph
construction. The vertex set has
\[
 |X|^T\,2^{|\mathcal B|}
 =|X|^T\,2^{2^{(n+1)T}}
\]
vertices. All factors except $|X|$ are constants for fixed $t,\eps$.
There are polynomially many endpoint and edge tuples, and all
word-generation sample spaces have constant size. Enumerate every test
outcome and add its probability to the corresponding edge weight.
The initial game weights, $t$, and $r$ are rational, so these weights
have polynomial bit complexity. Conditional edge probabilities also
have polynomial bit complexity. All alphabets and words can be enumerated
in constant time relative to the input size. The resulting graph is
therefore computable in deterministic polynomial time. It has a rational
edge distribution and satisfies the asserted, weaker $\eps$-loss bounds.
\end{proof}

\begin{proof}[Proof of \Cref{thm:main}]
Substituting $\rho=-t$ into \eqref{eq:gw-constant} gives
\begin{equation}\label{eq:gw-ratio}
 \GW=\min_{-1<t\le1}\frac{1+B(t)}{1+t},
\end{equation}
where the same formula defines $B(t)=2\arcsin(t)/\pi$ for negative $t$.
For $-1<t\le0$ the ratio is at least one, whereas for $0<t<1$ it is
less than one. At $t=1$ it equals one, and as $t\downarrow-1$ it
diverges. Thus its minimum occurs in $(0,1)$.

Fix $\GW<\alpha\le1$. By continuity and the density of the rationals,
choose a rational $t\in(0,1)$ such that
\[
 \alpha\frac{1+t}{2}>\frac{1+B(t)}{2}.
\]
Choose positive rational $\eps$ and $\epsilon'$ small enough that
$\eps<(t-B(t))/4$ and
\begin{equation}\label{eq:strict-final-gap}
 \alpha\left(\frac{1+t}{2}-\eps-\epsilon'\right)
 >\frac{1+B(t)}{2}+\eps+\epsilon'.
\end{equation}
Apply \Cref{thm:gap}. Delete loops while retaining the original weight
scale, so that the remaining weights sum to at most one. Apply
\Cref{lem:unweighted} with error $\epsilon'$, producing a simple
unweighted graph whose maximum cut, on the common scale $Q$, differs
from the original maximum cut weight by at most $\epsilon'$.
An $\alpha$-approximation algorithm on that graph returns, in the YES
case, a cut of scaled value at least the left side of
\eqref{eq:strict-final-gap}. In the NO case, no cut has scaled value
greater than its right side. A fixed rational threshold strictly between them therefore
distinguishes the NP-hard promise cases. This proves the theorem.
\end{proof}
